\documentclass[aspectratio=169]{beamer}
\usepackage{hyperref}
\usepackage{iftex}
\ifPDFTeX
  \usepackage[utf8]{inputenc}
  \usepackage[T1]{fontenc}
\else
  % XeTeX (tectonic, que es lo que usa `pixi run slides`) o LuaTeX: Unicode nativo.
  \usepackage{fontspec}
\fi
\usepackage[spanish]{babel}
\usepackage{csquotes}
\usepackage{tribhuvan}

% other packages
\usepackage{latexsym,amsmath,xcolor,multicol,booktabs,calligra}
\usepackage{graphicx,listings,textcomp}
\usepackage{tikz}
\usetikzlibrary{arrows.meta,positioning,shapes.geometric}
% Resaltado de sintaxis para listados de ficheros (pixi.toml, Dockerfile...),
% precomputado con pygmentize por generar.py para no depender de -shell-escape.
\usepackage{fvextra}
\makeatletter
\def\PY@reset{\let\PY@it=\relax \let\PY@bf=\relax
  \let\PY@ul=\relax \let\PY@tc=\relax
  \let\PY@bc=\relax \let\PY@ff=\relax}
\def\PY@tok#1{\csname PY@tok@#1\endcsname}
\def\PY@toks#1+{\ifx\relax#1\empty\else
  \PY@tok{#1}\expandafter\PY@toks\fi}
\def\PY@do#1{\PY@bc{\PY@tc{\PY@ul{\PY@it{\PY@bf{\PY@ff{#1}}}}}}}
\def\PY#1#2{\PY@reset\PY@toks#1+\relax+\PY@do{#2}}

\def\PY@color#1{\def\PY@tc##1{\textcolor[rgb]{#1}{##1}}}
\@namedef{PY@tok@w}{\PY@color{0.73,0.73,0.73}}
\@namedef{PY@tok@c}{\let\PY@it=\textit\PY@color{0.24,0.48,0.48}}
\@namedef{PY@tok@c1}{\let\PY@it=\textit\PY@color{0.24,0.48,0.48}}
\@namedef{PY@tok@c+c1}{\let\PY@it=\textit\PY@color{0.24,0.48,0.48}}
\@namedef{PY@tok@k}{\let\PY@bf=\textbf\PY@color{0.00,0.50,0.00}}
\@namedef{PY@tok@k+kc}{\let\PY@bf=\textbf\PY@color{0.00,0.50,0.00}}
\@namedef{PY@tok@o}{\PY@color{0.40,0.40,0.40}}
\@namedef{PY@tok@p}{\PY@color{0.20,0.20,0.20}}
\@namedef{PY@tok@n}{\PY@color{0.10,0.09,0.49}}
\@namedef{PY@tok@n+nt}{\let\PY@bf=\textbf\PY@color{0.00,0.50,0.00}}
\@namedef{PY@tok@n+nv}{\PY@color{0.10,0.09,0.49}}
\@namedef{PY@tok@l}{\PY@color{0.73,0.13,0.13}}
\@namedef{PY@tok@l+s+s2}{\PY@color{0.73,0.13,0.13}}
\@namedef{PY@tok@l+lScalar+lScalarPlain}{\PY@color{0.73,0.13,0.13}}
\@namedef{PY@tok@p+pIndicator}{\PY@color{0.40,0.40,0.40}}
\@namedef{PY@tok@err}{\PY@color{0.89,0.00,0.00}}
\@namedef{PY@tok@s}{\PY@color{0.73,0.13,0.13}}
\@namedef{PY@tok@s+se}{\PY@color{0.73,0.13,0.13}}
\@namedef{PY@tok@m}{\PY@color{0.40,0.40,0.40}}
\@namedef{PY@tok@nt}{\let\PY@bf=\textbf\PY@color{0.00,0.50,0.00}}

\def\PYZbs{\char`\\}\def\PYZus{\char`\_}
\def\PYZob{\char`\{}\def\PYZcb{\char`\}}
\def\PYZca{\char`\^}\def\PYZam{\char`\&}
\def\PYZlt{\char`\<}\def\PYZgt{\char`\>}
\def\PYZsh{\char`\#}\def\PYZpc{\char`\%}
\def\PYZdl{\char`\$}\def\PYZhy{\char`\-}
\def\PYZsq{\char`\'}\def\PYZdq{\char`\"}
\def\PYZti{\char`\~}\def\PYZat{@}
\def\PYZlb{[}\def\PYZrb{]}
\makeatother


\author[Francisco Martín Rico, Esteve Fernández, Juan Sebastián Cely, Francisco Miguel Moreno]{%
\begin{tabular}{cc}
Francisco Martín Rico & Esteve Fernández \\
\texttt{\footnotesize fmrico@gmail.com} & \texttt{\footnotesize esteve.fernandez@gmail.com} \\[1em]
Juan Sebastián Cely & Francisco Miguel Moreno \\
\texttt{\footnotesize juan.cely@urjc.es} & \texttt{\footnotesize franciscom.moreno@urjc.es}
\end{tabular}%
}
\title[Pixi y RoboStack]{Desde que uso Pixi estoy tan Pichi}
\subtitle{Workspaces ROS~2 declarativos y reproducibles con Pixi y RoboStack}
\institute{}

\date{}
\setbeamertemplate{date}{}
\setbeamertemplate{navigation symbols}{} % ocultar iconos de navegación

% defs
\definecolor{deepblue}{rgb}{0,0,0.5}
\definecolor{deepred}{rgb}{0.6,0,0}
\definecolor{deepgreen}{rgb}{0,0.5,0}
\definecolor{halfgray}{gray}{0.55}

\lstset{
    basicstyle=\ttfamily\small,
    keywordstyle=\bfseries\color{deepblue},
    emphstyle=\ttfamily\color{deepred},
    stringstyle=\color{deepgreen},
    numbers=none,
    breaklines=true,
    showstringspaces=false,
    rulesepcolor=\color{red!20!green!20!blue!20},
    frame=shadowbox,
    literate=
      {á}{{\'a}}1 {é}{{\'e}}1 {í}{{\'i}}1 {ó}{{\'o}}1 {ú}{{\'u}}1
      {Á}{{\'A}}1 {É}{{\'E}}1 {Í}{{\'I}}1 {Ó}{{\'O}}1 {Ú}{{\'U}}1
      {ñ}{{\~n}}1 {Ñ}{{\~N}}1
      {ü}{{\"u}}1 {Ü}{{\"U}}1
      {¿}{{\textquestiondown}}1 {¡}{{\textexclamdown}}1
}

% Frame auxiliar para preguntas/mini-retos interactivos al público
\newcommand{\preguntaframe}[1]{%
\begin{frame}
  \begin{center}
  \begin{beamercolorbox}[sep=1em,center]{block title}
  \Large \textbf{?} \; #1
  \end{beamercolorbox}
  \end{center}
\end{frame}%
}

% Frame auxiliar para marcar el inicio de un bloque práctico
\newcommand{\practicaframe}[2]{%
\begin{frame}
  \centering
  \vfill
  {\Large \textbf{PRÁCTICA}}\\[1em]
  {\large #1}\\[0.5em]
  {\color{halfgray} #2}
  \vfill
\end{frame}%
}

% Referencia al pie de una slide
\newcommand{\referencia}[1]{\vfill{\scriptsize\color{halfgray} #1}}

% Slide de separación: título grande y subtítulo, sin título de frame.
%   \separador{¿Por qué es tan simple?}{No tenéis que construir todo ROS vosotros.}
\newcommand{\separador}[2]{%
\begin{frame}
  \centering
  \vfill
  {\usebeamercolor[fg]{structure}\fontsize{34}{40}\selectfont\bfseries #1\par}
  \vspace{1.2em}
  {\Large\color{halfgray} #2\par}
  \vfill
\end{frame}%
}

% Inicio de sección: la entrada del índice y su slide de separación.
%   \seccion[Corto]{Título}{Subtítulo}
% Sustituye al índice automático que tribhuvan.sty pone al empezar cada sección.
\AtBeginSection[]{}
\newcommand{\seccion}[3][]{%
  \if\relax\detokenize{#1}\relax\section{#2}\else\section[#1]{#2}\fi
  \separador{#2}{#3}%
}

\begin{document}

%-------------------------------------------------
\begin{frame}
    \titlepage
    \begin{tikzpicture}[remember picture, overlay]
      \node[anchor=south east, inner sep=0pt] at ([xshift=-0.4cm, yshift=1.0cm]current page.south east)
        {\includegraphics[height=2.0cm]{figs/RosConES2026_small.png}};
      \node[anchor=south west, inner sep=0pt] (urjc) at ([xshift=1.0cm, yshift=1.0cm]current page.south west)
        {\includegraphics[width=2.1cm]{figs/logo_urjc.png}};
      \node[anchor=south west, inner sep=0pt] at ([yshift=0.3cm]urjc.north west)
        {\includegraphics[width=2.1cm]{figs/logo_irl.png}};
    \end{tikzpicture}
\end{frame}

\begin{frame}{Hoy}
\centering
\begin{tabular}{lll}
\toprule
\textbf{Duración} & \textbf{Bloque} & \textbf{Formato} \\
\midrule
10 min & Bienvenida y comprobación de la preparación & juntos \\
10 min & Robótica para todos & hablamos \\
25 min & Pixi en 25 minutos & hablamos \\
40 min & \textbf{Ejercicio 1:} tu primer workspace ROS~2 & escribís \\
15 min & Construir paquetes ROS con Pixi & hablamos \\
35 min & \textbf{Ejercicio 2:} construir y publicar paquetes & escribís \\
10 min & Colaboración, CI/CD y Docker & hablamos \\
20 min & \textbf{Ejercicio 3:} listo para tu equipo & escribís \\
20 min & \textbf{Ejercicio 4:} compilar un workspace real & escribís \\
\bottomrule
\end{tabular}
\referencia{Todo el material: \texttt{ejercicios/} (instrucciones paso a paso) y \texttt{soluciones/} (workspaces terminados) en el repositorio del taller.}
\end{frame}

%=================================================
\seccion{Bienvenida}{Que vuestro workspace ROS funcione en cualquier máquina.}
%=================================================

\begin{frame}{Antes de empezar\ldots}
\begin{itemize}
    \item ¿Cuántos habéis tenido un \lstinline|rosdep install| que falla distinto en vuestro portátil que en el del compañero?
    \item ¿Cuántos tenéis un \texttt{Dockerfile} de ROS~2 de más de 150 líneas solo para fijar versiones?
    \item ¿Cuántos tenéis instalado más de un \lstinline|/opt/ros/<distro>| en la misma máquina a la vez?
\end{itemize}
\vfill
\begin{block}{}
\centering No sois los únicos. Este taller va de eso.
\end{block}
\end{frame}

\begin{frame}[fragile]{¿Lo tenéis preparado?}
\begin{terminal}
curl -fsSL https://pixi.sh/install.sh | sh     # pixi --version >= 0.77
git clone <repositorio-del-taller>
cd workshop-pixi
pixi install --all --manifest-path soluciones/01-workspace-ros/pixi.toml
pixi run --manifest-path soluciones/01-workspace-ros/pixi.toml sim
\end{terminal}
\vfill
\centering ¿Aparece una tortuga? Estáis listos.
\vfill
\begin{alertblock}{¿Tenéis ROS instalado en el sistema?}
Comentad el \lstinline|source /opt/ros/...| de vuestro \lstinline|~/.bashrc| durante el taller: sus variables se cuelan en el entorno de Pixi y mezclan distros.
\end{alertblock}
\end{frame}

%=================================================
\seccion{Robótica para todos}{Por qué la forma de instalar ROS deja gente fuera.}
%=================================================

\begin{frame}{Entrar en ROS va atado a un SO y a una distro}
\centering
\begin{tabular}{ll}
\toprule
\textbf{Tienes} & \textbf{Tu camino} \\
\midrule
Ubuntu 24.04 & \lstinline|apt install|, el camino feliz \\
Ubuntu 22.04 & Solo Humble. ¿Quieres Jazzy? Reinstala el SO \\
Arch u otra Linux & Compilar desde fuente, o un contenedor \\
macOS & Sin camino oficial. VM o Docker \\
Windows & Soporte limitado, WSL2 o Docker \\
\bottomrule
\end{tabular}
\vfill
\begin{block}{Y además}
\lstinline|apt install ros-jazzy-desktop| hoy y dentro de 6 meses instala builds distintos: no hay lockfile. \texttt{rosdep} llama a \texttt{apt}, \texttt{pip}\ldots\ sin fijar versiones. Y el modelo underlay/overlay (\lstinline|source ... setup.bash|) es manual y frágil.
\end{block}
\end{frame}

\begin{frame}{``Usa Docker'' es la respuesta aceptada}
\begin{itemize}
    \item La instalación se vuelve posible, y el equipo comparte una imagen.
    \item Pero tu flujo de trabajo vive dentro de una caja: editor, depurador, USB, red, GPU, GUIs\ldots\ todo cuesta más.
    \item Docker aísla, pero \textbf{no fija el contenido}: un \lstinline|apt-get install| dentro del Dockerfile no es reproducible si no congelas todo a mano.
    \item Se monta una vez, reconstruir es lento, y nadie lo vuelve a tocar.
\end{itemize}
\vfill
\begin{alertblock}{Mensaje clave}
Docker sigue siendo válido para desplegar, pero no debería ser la \textbf{única} fuente de verdad de las dependencias del proyecto.
\end{alertblock}
\end{frame}

\begin{frame}[fragile]{Lo que queremos}
\vfill
\begin{terminal}
git clone tu-proyecto
cd tu-proyecto
pixi run application
\end{terminal}
\vfill
\centering \Large Cualquier máquina, tres comandos, y tu proyecto de robótica funciona.
\vfill
\normalsize ¿El equivalente a un \lstinline|package-lock.json| o un \lstinline|Cargo.lock|, \textbf{para ROS~2}?
\vfill
\end{frame}

\separador{¿Cómo llegamos ahí?}{Todas las máquinas · Sin distro · Declarativo · Reproducible · Rápido · Simple}

\begin{frame}{Cómo llegar (1/2)}
\begin{block}{Todas las máquinas}
Linux, macOS, Windows; x86 y arm64. Tu portátil, la CI y el robot pueden ser máquinas distintas.
\end{block}
\begin{block}{Sin distro}
Una distro de ROS es un proceso de release, no un sistema operativo. Jazzy en Ubuntu 22.04, Kilted en macOS, Humble en Windows\ldots\ y dos a la vez.
\end{block}
\begin{block}{Declarativo}
Tu entorno es un fichero en git: \lstinline|pixi.toml| dice lo que quieres. Vuelves a un estado bueno como vuelves atrás en el código.
\end{block}
\end{frame}

\begin{frame}{Cómo llegar (2/2)}
\begin{block}{Reproducible}
Un lockfile (\lstinline|pixi.lock|) en lugar de una imagen: cada paquete, versión y hash. El mismo entorno en el portátil, en la CI y en el robot.
\end{block}
\begin{block}{Rápido}
Minutos hasta tener el proyecto funcionando. Que lo único lento sea la conexión.
\end{block}
\begin{block}{Simple}
Si sabes usar \texttt{pip}, sabes usar Pixi: \lstinline|pixi init|, \lstinline|pixi add|, \lstinline|pixi run|. Sin \texttt{apt}, sin \texttt{rosdep}, sin \texttt{source}, sin \texttt{sudo}.
\end{block}
\end{frame}

\separador{¿Por qué es tan simple?}{No tenéis que construir vosotros todo ROS.}

\begin{frame}{Sobre hombros de gigantes: conda-forge y RoboStack}
\begin{center}
\resizebox{0.92\linewidth}{!}{%
\begin{tikzpicture}[
  caja/.style={draw, rounded corners=3pt, minimum width=4.6cm, minimum height=1.6cm, align=center, line width=0.6pt},
  fuente/.style={caja, draw=gray!60, fill=white},
  build/.style={caja, draw=orange!70!black, fill=yellow!20},
  canal/.style={caja, draw=blue!60, fill=blue!7},
  flecha/.style={-{Latex[length=3mm]}, line width=1pt},
  etiqueta/.style={above, font=\small},
  titulo/.style={font=\large\bfseries},
]
\node[fuente] (oss) {{\large\bfseries Proyectos open source}\\[2pt]Python · OpenCV · CMake};
\node[build, right=1.4cm of oss] (cfb) {{\large\bfseries Builds de conda-forge}\\[2pt]fuentes + recetas conda};
\node[canal, right=1.4cm of cfb] (cfc) {{\large\bfseries Canal conda-forge}\\[2pt]paquetes conda precompilados};
\node[fuente, below=2.4cm of oss] (ros) {{\large\bfseries Fuentes de paquetes ROS}\\[2pt]ROS Index / rosdistro};
\node[build, right=1.4cm of ros] (rsb) {{\large\bfseries Builds de RoboStack}\\[2pt]fuentes ROS + recetas conda};
\node[canal, right=1.4cm of rsb] (rsc) {{\large\bfseries Canales de RoboStack}\\[2pt]robostack-lyrical, -jazzy, \ldots};
\draw[flecha] (oss) -- node[etiqueta] {build} (cfb);
\draw[flecha] (cfb) -- node[etiqueta] {publish} (cfc);
\draw[flecha] (ros) -- node[etiqueta] {build} (rsb);
\draw[flecha] (rsb) -- node[etiqueta] {publish} (rsc);
\draw[flecha, draw=blue!60] (cfc.south) -- ++(0,-1.2) -| node[pos=0.25, below, align=center, font=\small]
  {{\color{blue!60}\bfseries Entorno de build}\\compiladores, herramientas y librerías} (rsb.north);
\end{tikzpicture}%
}
\end{center}
\vfill
\centering Pixi combina paquetes de \textbf{conda-forge} y de \textbf{RoboStack} en un mismo entorno.
\end{frame}

\begin{frame}{Un dato honesto}
\begin{alertblock}{RoboStack es hoy Tier 3}
En la clasificación oficial de plataformas de ROS: soporte \textbf{comunitario}, no de Open Robotics.
\end{alertblock}
\vfill
\begin{itemize}
    \item Con adopción creciente, y anunciado como parte de la estrategia de accesibilidad de Open Robotics.
    \item Mantenerlo es trabajo duro: Tobias Fischer, Silvio Traversaro, Daisuke Nishimatsu y muchos más.
    \item ¿Falta un paquete en RoboStack? Es un pull request a \texttt{vinca.yaml}.
\end{itemize}
\end{frame}

%=================================================
\seccion[Pixi]{Pixi en 25 minutos}{Lo básico que necesitáis para el Ejercicio 1.}
%=================================================

\begin{frame}{¿Qué es Pixi?}
\centering
\begin{tabular}{ll}
\toprule
\textbf{Gestor de paquetes} & \textbf{Qué instala} \\
\midrule
\texttt{apt} & Paquetes Debian \\
\texttt{brew} & Paquetes Homebrew \\
\texttt{pip} & Paquetes Python de PyPI \\
\texttt{pixi} & \textbf{paquetes conda + paquetes Python de PyPI} \\
\bottomrule
\end{tabular}
\vfill
\begin{itemize}
    \item Escrito en Rust sobre \textbf{Rattler} (reimplementación del ecosistema conda), por el equipo de \textbf{prefix.dev}, autores de \textbf{mamba}.
    \item Sin entorno \texttt{base}: un binario estático. Todo vive en la carpeta del proyecto, con \lstinline|pixi.toml| + \lstinline|pixi.lock|.
\end{itemize}
\end{frame}

\preguntaframe{Con lo que sabéis de \texttt{pip} o \texttt{npm}, ¿cómo creéis que se instala ROS~2 con Pixi?}

\begin{frame}{Instalar ROS~2}
\centering
\scriptsize
\begin{tabular}{ll}
\toprule
\textbf{Instalación oficial en Ubuntu} & \textbf{Con Pixi + RoboStack} \\
\midrule
Ubuntu 24.04 y un locale UTF-8 & Pixi en Linux, macOS o Windows \\
Universe + \texttt{ros2-apt-source} & \texttt{pixi init -c robostack-lyrical -c conda-forge} \\
\texttt{sudo apt update \&\& sudo apt upgrade} & Sin tocar el sistema \\
\texttt{sudo apt install ros-lyrical-desktop} & \texttt{pixi add ros-lyrical-desktop} \\
\texttt{source /opt/ros/lyrical/setup.bash} & \texttt{pixi shell} \\
\texttt{ros2 run demo\_nodes\_cpp talker} & \texttt{ros2 run demo\_nodes\_cpp talker} \\
\bottomrule
\end{tabular}
\end{frame}

\begin{frame}[fragile]{El \texttt{pixi.toml}}
%%CODE toml
[workspace]
# Para qué máquinas se resuelve (y se genera el lockfile)
platforms = ["linux-64", "osx-arm64", "win-64"]
# De dónde salen los paquetes: RoboStack para ROS, conda-forge para el resto
channels = ["robostack-lyrical", "conda-forge"]

[dependencies]
ros-lyrical-desktop = "*"

[tasks]
start = "ros2 launch my_package my_launch_file.launch.py"
%%END
\vfill
Vive en la \textbf{raíz} de tu workspace. Además: \lstinline|[activation]|, \lstinline|[environments]|, \lstinline|[target.<plataforma>]|, \lstinline|[pypi-dependencies]|\ldots
\end{frame}

\begin{frame}{El comando \texttt{pixi}}
\centering
\begin{tabular}{ll}
\toprule
\textbf{Comando} & \textbf{Qué hace} \\
\midrule
\lstinline|pixi init| & Crea un workspace \\
\lstinline|pixi add ros-lyrical-ros-base| & Añade la dependencia, resuelve e instala \\
\lstinline|pixi run ros2 topic list| & Ejecuta un comando dentro del entorno \\
\lstinline|pixi shell| & Abre un shell dentro del entorno \\
\lstinline|pixi install| & Instala el entorno desde el lockfile \\
\lstinline|pixi list| / \lstinline|pixi tree| & Qué hay instalado, y por qué \\
\lstinline|pixi global install gh| & Herramientas CLI globales, aisladas \\
\bottomrule
\end{tabular}
\referencia{\lstinline|pixi <comando> --help| · chuleta en \texttt{referencia/chuleta.md}}
\end{frame}

\begin{frame}[fragile]{Tasks}
\begin{columns}[T]
\column{0.55\linewidth}
%%CODE toml
[tasks]
sim = "ros2 run turtlesim turtlesim_node"

[tasks.build]
cmd = "colcon build"
inputs = ["src/turtle_dancer"]

[tasks.dance]
cmd = "ros2 run turtle_dancer dance"
depends-on = ["build"]
%%END
\column{0.42\linewidth}
\begin{terminal}
# En todas las plataformas
pixi run sim

# Compila antes
pixi run dance

# Cualquier comando
pixi run ros2 topic list
\end{terminal}
\end{columns}
\vfill
\begin{itemize}
    \item Se ejecutan en un shell multiplataforma propio (\texttt{deno\_task\_shell}), dentro del entorno.
    \item \lstinline|depends-on| encadena; \lstinline|inputs|/\lstinline|outputs| cachean: no se repite si nada cambió.
\end{itemize}
\end{frame}

\begin{frame}[fragile]{\texttt{pixi shell} y la activación}
\begin{columns}[T]
\column{0.45\linewidth}
\begin{terminal}
pixi shell
ros2 topic list
ros2 run turtlesim turtlesim_node
exit
\end{terminal}
\column{0.5\linewidth}
%%CODE toml
[activation]
scripts = ["install/setup.sh"]

[activation.env]
ROS_DOMAIN_ID = "42"
%%END
\end{columns}
\vfill
\begin{itemize}
    \item \lstinline|pixi shell| te mete en el entorno, como hacía \lstinline|source install/setup.bash|.
    \item \lstinline|pixi run| y \lstinline|pixi shell| activan el entorno primero, incluidos tus scripts.
    \item \lstinline|pixi shell-hook| imprime la activación como un script de shell normal (lo usaremos en Docker).
\end{itemize}
\end{frame}

\begin{frame}{El lockfile: el porqué de la reproducibilidad}
\begin{itemize}
    \item \lstinline|pixi.lock| fija versión exacta, build y \textbf{sha256} de cada paquete (conda y PyPI), \textbf{para cada plataforma declarada}.
    \item Se actualiza solo con \lstinline|pixi add| / \lstinline|pixi update|.
    \item \lstinline|pixi install --locked| $\rightarrow$ falla si el lockfile no está sincronizado (la CI).
    \item \lstinline|pixi lock --check| $\rightarrow$ la barrera barata en cada pull request.
\end{itemize}
\vfill
\begin{block}{Se commitea a git}
Todo el equipo, la CI y el robot instalan exactamente los mismos binarios. Es más fácil tirarlo que reconstruirlo cuando algo se actualizó sin que te enterases\ldots\ justo antes de la demo.
\end{block}
\end{frame}

\begin{frame}[fragile]{Compilar tu paquete con colcon}
Con ROS instalado y \lstinline|src/turtle_dancer/| en tu workspace:
\begin{terminal}
pixi add ros-dev-tools        # colcon, CMake y compiladores
pixi run colcon build
pixi workspace activation scripts add --target unix install/setup.sh
pixi run ros2 run turtle_dancer dance
\end{terminal}
\vfill
\begin{itemize}
    \item Igual que siempre, pero dentro del entorno: aparecen \lstinline|build/|, \lstinline|install/| y \lstinline|log/|.
    \item La activación carga el overlay por ti. Nada de \texttt{setuptools} ni \texttt{empy} fijados a mano.
\end{itemize}
\end{frame}

\begin{frame}[fragile]{Dos distros, un workspace}
%%CODE toml
[environments.default]
channels = ["robostack-lyrical", "conda-forge"]
[environments.default.dependencies]
ros-lyrical-ros-base = ">=0.13"

[environments.jazzy]
channels = ["robostack-jazzy", "conda-forge"]
[environments.jazzy.dependencies]
ros-jazzy-ros-base = "*"
%%END
\vfill
Cada entorno tiene sus canales y paquetes; las tasks se comparten.
\lstinline|pixi run sim| para Lyrical, \lstinline|pixi run -e jazzy sim| para Jazzy.
\vfill
\small Imposible o muy incómodo con \texttt{apt}. Con Pixi, un entorno más.
\end{frame}

\begin{frame}[fragile]{Resolver para máquinas que no tienes}
%%CODE toml
[workspace]
platforms = [
  "linux-64", "osx-arm64", "win-64",
  { name = "cuda-linux-64", platform = "linux-64", cuda = "12" },
  { name = "jetson", platform = "linux-aarch64", cuda = "13" },
]
%%END
\begin{terminal}
pixi lock                          # resuelve todas las plataformas
pixi list --platform jetson        # lo que tendría la Jetson
\end{terminal}
\vfill
\centering Resolver es barato y portable; instalar es lo caro, y se hace donde toca.
\end{frame}

\begin{frame}[fragile]{Paquetes virtuales y CUDA}
\begin{columns}[T]
\column{0.34\linewidth}
\begin{terminal}
$ pixi info
Virtual packages:
  __unix=0=0
  __linux=6.8.0=0
  __glibc=2.39=0
  __cuda=12.4=0
  __archspec=1=x86_64
\end{terminal}
\column{0.63\linewidth}
%%CODE toml
[dependencies]
# GPU donde hay CUDA, CPU en el resto
pytorch-gpu = { version = ">=2.5", when = "__cuda" }
pytorch = ">=2.5"
%%END
\end{columns}
\vfill
\begin{itemize}
    \item No son paquetes que instalas: son hechos sobre la máquina. El solver los usa.
    \item En una máquina que no tienes, los declaras tú (\lstinline|cuda = "12"|).
    \item El driver viene del SO; todo lo demás (toolkit, cuDNN, PyTorch) de conda-forge.
\end{itemize}
\end{frame}

\begin{frame}[fragile]{Un workspace nuevo en cuatro líneas}
\vfill
\begin{terminal}
pixi init ros-workspace -c robostack-lyrical -c conda-forge
cd ros-workspace
pixi add ros-lyrical-ros-base ros-lyrical-turtlesim
pixi run ros2 run turtlesim turtlesim_node
\end{terminal}
\vfill
\end{frame}

\practicaframe{Ejercicio 1: tu primer workspace ROS~2}{40 minutos · \texttt{cd ejercicios/01-workspace-ros}}

\begin{frame}{Ejercicio 1: tu primer workspace ROS~2}
\begin{enumerate}
    \item \lstinline|pixi init|, canal \texttt{robostack-lyrical}, \lstinline|ros-lyrical-ros-base| + \lstinline|turtlesim|.
    \item Tasks \texttt{sim}, \texttt{teleop}, \texttt{topics}.
    \item Compilar \texttt{turtle\_dancer} con colcon y cargar el overlay con \lstinline|[activation]|.
    \item Lyrical y Jazzy en el mismo workspace.
    \item Un cerebro con PyTorch (\texttt{turtle\_brain}).
    \item Plataformas: CPU, GPU (\texttt{cuda-linux-64}) y Jetson, resueltas desde tu portátil.
    \item (Opcional) Ejecutarlo en una GPU de verdad.
\end{enumerate}
\vfill
\small Instrucciones: \texttt{ejercicios/01-workspace-ros/README.md} · ¿Atascado? \texttt{soluciones/01-workspace-ros/} o levanta la mano.
\end{frame}

%=================================================
\seccion[Paquetes]{Construir paquetes ROS con Pixi}{Sin colcon, sin overlay, sin \texttt{source}. Y sin montarse una buildfarm.}
%=================================================

\begin{frame}{Antes: tu propia buildfarm}
\begin{center}
\resizebox{\linewidth}{!}{%
\begin{tikzpicture}[node distance=0.7cm, every node/.style={font=\footnotesize}]
\node[draw, rounded corners, fill=green!15, minimum width=2.4cm, minimum height=0.8cm] (a) {\texttt{recipe.yaml}};
\node[draw, rounded corners, fill=green!20, minimum width=2.4cm, minimum height=0.8cm, right=of a] (b) {rattler-build};
\node[draw, rounded corners, fill=green!20, minimum width=2.4cm, minimum height=0.8cm, right=of b] (c) {rattler-index};
\node[draw, rounded corners, fill=orange!15, minimum width=2.4cm, minimum height=0.8cm, right=of c] (d) {\texttt{upload prefix}};
\node[draw, rounded corners, fill=blue!12, minimum width=2.4cm, minimum height=0.8cm, right=of d] (e) {\texttt{pixi.toml}};
\draw[-Latex, thick] (a) -- (b);
\draw[-Latex, thick] (b) -- (c);
\draw[-Latex, thick] (c) -- (d);
\draw[-Latex, thick] (d) -- (e);
\end{tikzpicture}%
}
\end{center}
\begin{itemize}
    \item Una receta a mano por paquete, duplicando lo que ya dice el \lstinline|package.xml|.
    \item Para ROS completo: fork de RoboStack, \texttt{vinca}, CI propia\ldots
\end{itemize}
\vfill
\begin{center}
\resizebox{0.75\linewidth}{!}{%
\begin{tikzpicture}[node distance=0.7cm, every node/.style={font=\footnotesize}]
\node[draw, rounded corners, fill=blue!12, minimum width=3.6cm, minimum height=0.8cm] (a) {\texttt{package.xml} + \texttt{pixi.toml}};
\node[draw, rounded corners, fill=green!20, minimum width=2.4cm, minimum height=0.8cm, right=of a] (b) {\texttt{pixi publish}};
\node[draw, rounded corners, fill=orange!15, minimum width=2.4cm, minimum height=0.8cm, right=of b] (c) {tu canal};
\draw[-Latex, thick] (a) -- (b);
\draw[-Latex, thick] (b) -- (c);
\end{tikzpicture}%
}
\end{center}
\centering \textbf{Ahora}: dos líneas junto al \lstinline|package.xml|, y un comando.
\end{frame}

\begin{frame}[fragile]{Dos papeles para un \texttt{pixi.toml}}
\centering
\begin{tabular}{lll}
\toprule
\textbf{Workspace} & \lstinline|[workspace]| & lo que ejecutas con \lstinline|pixi run| \\
\textbf{Paquete} & \lstinline|[package]| & lo que Pixi construye \\
\bottomrule
\end{tabular}
\vfill
\raggedright
Todavía en \emph{preview}: puede cambiar algo de sintaxis.
\begin{terminal}
pixi workspace preview add pixi-build
pixi workspace dependencies add "pixi-build-ros >=0.7.5"
\end{terminal}
\end{frame}

\begin{frame}[fragile]{Lo que escribes}
\begin{columns}[T]
\column{0.48\linewidth}
\small \texttt{pixi.toml} (workspace)
%%CODE toml
[workspace]
preview = ["pixi-build"]

[workspace.dependencies]
pixi-build-ros = ">=0.7.5"

[dependencies.ros-lyrical-turtle-dancer]
path = "src/turtle_dancer"
%%END
\column{0.48\linewidth}
\small \texttt{src/turtle\_dancer/pixi.toml} (paquete)
%%CODE toml
[package.build.backend]
name = "pixi-build-ros"
workspace = true
%%END
\end{columns}
\vfill
\begin{itemize}
    \item Nombre de la dependencia: el del \lstinline|package.xml| con el prefijo de la distro y guiones.
    \item \lstinline|CMakeLists.txt| y \lstinline|setup.py| no se tocan: el backend ejecuta tu build \texttt{ament}.
\end{itemize}
\end{frame}

\begin{frame}[fragile]{\texttt{pixi-build-ros} lee tu \texttt{package.xml}}
\begin{columns}[T]
\column{0.48\linewidth}
\small Lo que ya tienes
%%CODE xml tiny
<package format="3">
  <name>turtle_dancer</name>
  <version>0.1.0</version>
  <buildtool_depend>ament_cmake</buildtool_depend>
  <depend>rclcpp</depend>
  <depend>geometry_msgs</depend>
  <exec_depend>turtlesim</exec_depend>
  <export>
    <build_type>ament_cmake</build_type>
  </export>
</package>
%%END
\column{0.48\linewidth}
\small Lo que \textbf{no} tienes que escribir
%%CODE toml tiny
[package]
name = "ros-lyrical-turtle-dancer"
version = "0.1.0"
[package.build-dependencies]
ros-lyrical-ament-cmake = "*"
ros-lyrical-rclcpp = "*"
ros-lyrical-geometry-msgs = "*"
[package.run-dependencies]
ros-lyrical-rclcpp = "*"
ros-lyrical-geometry-msgs = "*"
ros-lyrical-turtlesim = "*"
%%END
\end{columns}
\vfill
\centering \small El \lstinline|package.xml| sigue siendo la única fuente de verdad. La distro la elige el canal del workspace.
\end{frame}

\begin{frame}{El bucle editar-ejecutar}
\begin{itemize}
    \item \lstinline|pixi install| construye el paquete en un entorno de build aislado y lo instala en tu entorno \textbf{como un paquete conda}, al lado de \texttt{turtlesim}.
    \item \lstinline|pixi run dance| comprueba antes las entradas (fuentes, \lstinline|package.xml|, \lstinline|CMakeLists.txt|, launch, mensajes\ldots) y reconstruye si algo cambió.
    \item Fuera: la task \texttt{build}, \lstinline|ros-dev-tools|, la activación, \lstinline|build/ install/ log/|.
\end{itemize}
\vfill
\begin{block}{}
\centering Editar $\rightarrow$ \lstinline|pixi run| $\rightarrow$ funciona. Desde una terminal nueva, sin \texttt{source} de nada.
\end{block}
\end{frame}

\begin{frame}{Backends de build}
Pixi no sabe compilar tu código. Un \textbf{backend} sí.
\vfill
\centering
\begin{tabular}{lll}
\toprule
\textbf{Backend} & \textbf{Construye} & \textbf{Lee} \\
\midrule
\texttt{pixi-build-ros} & Paquetes ROS & \lstinline|package.xml| \\
\texttt{pixi-build-cmake} & C y C++ & \lstinline|CMakeLists.txt| \\
\texttt{pixi-build-python} & Python & \lstinline|pyproject.toml| \\
\texttt{pixi-build-rust} & Rust & \lstinline|Cargo.toml| \\
\texttt{pixi-build-rattler-build} & Recetas conda & \lstinline|recipe.yaml| \\
\bottomrule
\end{tabular}
\vfill
\raggedright \small ¿Necesitas control total (parches, pasos de build propios)? Tu \lstinline|recipe.yaml| de \texttt{rattler-build} sigue valiendo: \lstinline|[package.build.config] recipe = "recipe.yaml"|.
\end{frame}

\begin{frame}[fragile]{Publicar}
Cada paquete decide si se publica:
%%CODE toml
# src/turtle_dancer/pixi.toml
[package]
publish = true
%%END
\begin{terminal}
# Canal local, como en el Ejercicio 2
pixi publish --target-channel output

# Tu canal de prefix.dev (cuenta gratuita, permiso de escritura)
pixi auth login prefix.dev
pixi publish --target-channel https://prefix.dev/<tu-canal>
\end{terminal}
\vfill
\small Construye en orden de dependencias todos los paquetes con \lstinline|publish = true| y crea el índice del canal. Para otras plataformas: \lstinline|--target-platform|, o una matriz en la CI.
\end{frame}

\begin{frame}[fragile]{Depender de un paquete}
\begin{columns}[T]
\column{0.5\linewidth}
%%CODE toml
[workspace]
channels = [
  "https://prefix.dev/<tu-canal>",
  "https://prefix.dev/robostack-lyrical",
  "conda-forge",
]

[dependencies]
# Ya construido, desde un canal
ros-lyrical-turtle-dancer = "*"
%%END
\column{0.47\linewidth}
%%CODE toml
# Desde el código, local
[dependencies.ros-lyrical-turtle-dancer]
path = "src/turtle_dancer"

# Desde git
[dependencies.brain]
git = "https://github.com/user/brain"
branch = "main"
%%END
\end{columns}
\vfill
\small El orden de los canales importa: tu canal $\rightarrow$ RoboStack $\rightarrow$ conda-forge. El mismo patrón que la buildfarm, sin buildfarm.
\end{frame}

\practicaframe{Ejercicio 2: construir y publicar paquetes ROS}{35 minutos · \texttt{cd ejercicios/02-paquetes-ros}}

\begin{frame}{Ejercicio 2: construir y publicar paquetes ROS}
\begin{enumerate}
    \item Activar \texttt{pixi-build} y añadir el backend al workspace.
    \item Escribir \lstinline|src/turtle_dancer/pixi.toml|: dos líneas.
    \item Añadir la dependencia \texttt{path}, \lstinline|pixi install|.
    \item Borrar lo de colcon, \lstinline|pixi run dance|.
    \item Cambiar \texttt{dance.cpp}, volver a ejecutar.
    \item Lo mismo con \lstinline|src/turtle_choreographer| (Python).
    \item Publicar los dos paquetes en un canal (local o prefix.dev).
    \item Instalarlos desde otro workspace, sin código fuente.
\end{enumerate}
\vfill
\small Instrucciones: \texttt{ejercicios/02-paquetes-ros/README.md}
\end{frame}

%=================================================
\seccion[Colaboración]{Colaboración, CI/CD y Docker}{De «funciona en mi máquina» a «funciona en la del equipo».}
%=================================================

\begin{frame}[fragile]{CI en dos líneas}
%%CODE yaml
jobs:
  test:
    runs-on: ubuntu-latest
    steps:
      - uses: actions/checkout@v7
      - uses: prefix-dev/setup-pixi@v0.10.2
        with:
          locked: true
      - run: pixi run test
%%END
\begin{itemize}
    \item Instala Pixi y el entorno desde el lockfile; ejecuta tu task.
    \item Las mismas tasks que en tu portátil: no hay un segundo setup en YAML que se desincronice.
    \item Caché gratis, con el hash de \lstinline|pixi.lock| como clave.
\end{itemize}
\end{frame}

\begin{frame}[fragile]{Docker: la etapa de build}
%%CODE docker
FROM ghcr.io/prefix-dev/pixi:0.80.0-noble AS build
WORKDIR /app
COPY pixi.toml pixi.lock ./
RUN pixi install --locked
RUN pixi shell-hook -s bash > /shell-hook.sh \
    && echo 'exec "$@"' >> /shell-hook.sh
%%END
\vfill
\begin{itemize}
    \item Parte de la imagen de Pixi y copia el manifiesto y el lockfile.
    \item \lstinline|--locked|: falla si el lockfile no cuadra; la imagen no puede contener algo que el equipo no probó.
    \item Escribe la activación en un script, para la siguiente etapa.
\end{itemize}
\end{frame}

\begin{frame}[fragile]{Docker: la etapa final}
%%CODE docker
FROM ubuntu:26.04 AS runtime
COPY --from=build /app/.pixi/envs/default /app/.pixi/envs/default
COPY --from=build /shell-hook.sh /shell-hook.sh
ENTRYPOINT ["/bin/bash", "/shell-hook.sh"]
CMD ["ros2", "run", "demo_nodes_cpp", "talker"]
%%END
\vfill
\begin{itemize}
    \item Una imagen base normal, sin el binario de \texttt{pixi}.
    \item Recibe el entorno terminado y el script de activación.
    \item Portátil, CI e imagen instalan desde el mismo \lstinline|pixi.lock|.
\end{itemize}
\end{frame}

\begin{frame}[fragile]{\texttt{pixi-pack}: un entorno en un fichero}
\begin{columns}[T]
\column{0.45\linewidth}
\small En tu portátil
\begin{terminal}
pixi-pack \
  --platform linux-aarch64 \
  pixi.toml
\end{terminal}
\column{0.5\linewidth}
\small En el robot
\begin{terminal}
pixi-unpack environment.tar
source activate.sh
ros2 run demo_nodes_cpp talker
\end{terminal}
\end{columns}
\vfill
\begin{itemize}
    \item Para despliegues sin red, con 4G poco fiable o sin Docker.
    \item En el destino no hace falta Pixi, ni conda, ni conexión.
    \item \lstinline|--create-executable|: un único fichero autoextraíble.
\end{itemize}
\end{frame}

\practicaframe{Ejercicio 3: listo para tu equipo}{20 minutos · \texttt{cd ejercicios/03-colaboracion}}

\begin{frame}{Ejercicio 3: listo para tu equipo}
\begin{enumerate}
    \item \lstinline|pixi global install gh| y subir el workspace a GitHub.
    \item CI con \texttt{setup-pixi}: verla en verde.
    \item Completar el \texttt{Dockerfile}, construir la imagen y ejecutarla: sin Pixi dentro.
    \item Empaquetar el entorno, desempaquetarlo y ejecutar sin Pixi.
\end{enumerate}
\vfill
\small Instrucciones: \texttt{ejercicios/03-colaboracion/README.md}
\end{frame}


%=================================================
\seccion[Workspace real]{Compilar un workspace real}{Repositorios de otros, sin \texttt{rosdep}, y un Kobuki en una casa.}
%=================================================

\begin{frame}{La vida real: un \texttt{src/} con repositorios de otros}
\begin{itemize}
    \item Lo de siempre: clonar, \lstinline|rosdep install|, \lstinline|colcon build|, \lstinline|source|.
    \item \texttt{rosdep} no sirve dentro de Pixi: llama a \texttt{apt} y \texttt{pip} del sistema, con \texttt{sudo} y sin fijar versiones.
    \item Lo que hace \texttt{rosdep} es traducir cada dependencia del \lstinline|package.xml| a un paquete. Con RoboStack la traducción es directa:
\end{itemize}
\vfill
\centering
\scriptsize
\begin{tabular}{lll}
\toprule
\textbf{En el \texttt{package.xml}} & \textbf{Con \texttt{rosdep}} & \textbf{Con Pixi} \\
\midrule
\texttt{ros\_gz\_sim} & \texttt{apt install ros-lyrical-ros-gz-sim} & \texttt{pixi add ros-lyrical-ros-gz-sim} \\
\texttt{xacro} & \texttt{apt install ros-lyrical-xacro} & \texttt{pixi add ros-lyrical-xacro} \\
\bottomrule
\end{tabular}
\vfill
\raggedright\small Y queda escrito en el \lstinline|pixi.toml| y fijado en el \lstinline|pixi.lock|. Para workspaces grandes, \texttt{pixi-ros} hace la traducción por ti.
\end{frame}

\begin{frame}[fragile]{Los repositorios, en un fichero: \texttt{workspace.repos}}
%%CODE yaml
repositories:
  easynav_playgrounds:
    type: git
    url: https://github.com/EasyNavigation/easynav_playgrounds.git
    version: pichi-workshop      # solo el Kobuki y la casa
  teleop_twist_keyboard:
    type: git
    url: https://github.com/ros2/teleop_twist_keyboard.git
    version: rolling
%%END
\begin{terminal}
vcs import --input workspace.repos src     # vcstool viene en ros-dev-tools
\end{terminal}
\vfill
\small El formato estándar de ROS para describir un workspace: va a git junto al \lstinline|pixi.toml| y el \lstinline|pixi.lock|.
\end{frame}

\begin{frame}[fragile]{El entorno y las tasks: \texttt{pixi.toml}}
\begin{columns}[T]
\column{0.4\linewidth}
%%CODE toml
[dependencies]
ros-lyrical-ros-base = ">=0.13"
ros-dev-tools = ">=1.0"
ros-lyrical-ros-gz-sim = "*"
ros-lyrical-ros-gz-bridge = "*"
ros-lyrical-gz-sim-vendor = "*"
ros-lyrical-xacro = "*"
# ...

[target.unix.activation]
scripts = ["install/setup.sh"]
%%END
\column{0.57\linewidth}
%%CODE toml
[tasks.importar]
cmd = "mkdir -p src && vcs import --input workspace.repos src"

[tasks.build]
cmd = "colcon build"
depends-on = ["importar"]

[tasks.sim]
cmd = "ros2 launch easynav_playground_kobuki_worlds gazebo_sim.launch.yaml"
depends-on = ["build"]
%%END
\end{columns}
\vfill
\centering \lstinline|git clone| · \lstinline|cd| · \lstinline|pixi run sim|: clona, compila y lanza Gazebo con el Kobuki.
\end{frame}

\practicaframe{Ejercicio 4: compilar un workspace real}{20 minutos · \texttt{cd ejercicios/04-workspace-real}}

\begin{frame}{Ejercicio 4: compilar un workspace real}
\begin{enumerate}
    \item Workspace con \texttt{robostack-lyrical}, \lstinline|ros-lyrical-ros-base| y \lstinline|ros-dev-tools|.
    \item \lstinline|workspace.repos| con \texttt{easynav\_playgrounds} (rama \texttt{pichi-workshop}) y \texttt{teleop\_twist\_keyboard}; \lstinline|vcs import|.
    \item Leer los \lstinline|package.xml| y añadir sus dependencias con \lstinline|pixi add|: vuestro \lstinline|rosdep install|.
    \item \lstinline|colcon build| y la activación del overlay, como en el Ejercicio 1.
    \item \lstinline|pixi run sim|: Gazebo con el Kobuki en la casa.
    \item \lstinline|pixi run teleop| en otra terminal: pasead el Kobuki.
\end{enumerate}
\vfill
\small Instrucciones: \texttt{ejercicios/04-workspace-real/README.md} · Gazebo necesita OpenGL: en VMs o WSL2, \lstinline|LIBGL_ALWAYS_SOFTWARE=1|.
\end{frame}

%=================================================
\seccion{Cierre}{Lo que os lleváis, y dónde seguir.}
%=================================================

\begin{frame}{Lo que os lleváis}
\begin{itemize}
    \item Un workspace ROS~2 en un \lstinline|pixi.toml| + \lstinline|pixi.lock|: \lstinline|git clone|, \lstinline|cd|, \lstinline|pixi run|.
    \item Varias distros, varias plataformas y GPU/Jetson resueltas desde vuestro portátil.
    \item Paquetes ROS construidos desde su \lstinline|package.xml| y publicados con \lstinline|pixi publish|.
    \item La misma definición en el portátil, en la CI, en Docker y en el robot.
\end{itemize}
\end{frame}

\begin{frame}{Con vuestro propio workspace}
\begin{enumerate}
    \item \lstinline|pixi-ros init --distro <distro>| en la raíz: revisad los \texttt{NOT FOUND}.
    \item \lstinline|pixi run build| con colcon, como en el Ejercicio 1.
    \item Paquete a paquete (empezando por las hojas), un \lstinline|pixi.toml| de dos líneas y una dependencia \texttt{path}, como en el Ejercicio 2.
    \item \lstinline|publish = true| y \lstinline|pixi publish| desde la CI: vuestra ``buildfarm''.
    \item \lstinline|setup-pixi| en la CI y \lstinline|pixi-pack| o Docker para el robot, como en el Ejercicio 3.
\end{enumerate}
\vfill
\small Si una dependencia del \lstinline|package.xml| no está en RoboStack con ese nombre: \texttt{extra-package-mappings} en la configuración de \texttt{pixi-build-ros}.
\end{frame}

\begin{frame}{Recursos}
\small
\begin{itemize}
    \item Pixi: \url{https://pixi.prefix.dev} — ROS: \url{https://pixi.prefix.dev/latest/build/ros/}
    \item RoboStack: \url{https://robostack.github.io/} — \url{https://github.com/RoboStack}
    \item pixi-ros: \url{https://github.com/ruben-arts/pixi-ros}
    \item Discord de Pixi, canal \texttt{\#pixi}: \url{https://discord.gg/kKV8ZxyzY4}
    \item Paper original: Fischer et al., \emph{``A RoboStack Tutorial''}, IEEE RAM 2021 (arXiv:2104.12910)
\end{itemize}
\end{frame}

\begin{frame}{Agradecimientos}
\begin{block}{Basado en el taller de prefix.dev en ROSCon 2026}
\emph{Declarative ROS workspaces with Pixi and RoboStack}, de Ruben Arts, Wolf Vollprecht y Bas Zalmstra (BSD-3-Clause).\\[0.5em]
\small\url{https://prefix-dev.github.io/roscon-2026-declarative-ros-workspaces-with-pixi/}
\end{block}
\vfill
Los paquetes \texttt{turtle\_dancer}, \texttt{turtle\_choreographer} y \texttt{turtle\_brain} de los ejercicios vienen de ese taller.
\vfill
Y a la comunidad de conda-forge y RoboStack, sin la que nada de esto funcionaría.
\end{frame}

\begin{frame}
  \centering
  \Large Preguntas
\end{frame}

\end{document}
